\documentclass[11pt]{article}
\usepackage[a4paper,margin=27mm]{geometry}
\usepackage[T1]{fontenc}
\usepackage[utf8]{inputenc}
\usepackage{lmodern,amsmath,amssymb,amsthm,mathtools,microtype,booktabs}
\usepackage{hyperref}
\hypersetup{colorlinks=true,linkcolor=blue,citecolor=blue,urlcolor=blue}
\makeatletter
\renewcommand{\@seccntformat}[1]{\csname the#1\endcsname.\quad}
\makeatother
\newtheorem{theorem}{Theorem}
\newtheorem{proposition}[theorem]{Proposition}
\newtheorem{lemma}[theorem]{Lemma}
\newtheorem{corollary}[theorem]{Corollary}
\theoremstyle{remark}
\newtheorem{remark}[theorem]{Remark}
\DeclareMathOperator{\Hess}{Hess}
\DeclareMathOperator{\diag}{diag}
\DeclareMathOperator{\rank}{rank}
\DeclareMathOperator{\tr}{tr}
\newcommand{\R}{\mathbb R}
\newcommand{\one}{\mathbf 1}
\newcommand{\bnu}{\boldsymbol{\nu}}
\newcommand{\dd}{\partial}
\title{Rank-two obstructions to Lorentzian\\graph partition functions}
\author{Lluis Eriksson\\ \small Independent researcher, Sweden}
\date{17 October 2026\\ \small Research draft}
\begin{document}
\maketitle
\begin{abstract}
We give rank-two antiferromagnetic obstructions to the Lorentzian
property of graph partition functions. A connected noncomplete source
of maximum degree $\Delta$ has a strictly positive rational obstruction
on $\Delta+1$ spins whenever no two maximum-degree vertices have equal
closed neighbourhoods. Consequently, the connected-source
characterization holds for graphs without true twins, graphs with a
unique maximum-degree vertex, and all connected edge-regular graphs,
including strongly regular graphs. The proof uses a small-parameter
limit of explicit conditioned-set Hessian kernels.
A different endpoint yields a neighbourhood-independence-gap criterion.
One fixed nonnegative four-spin matrix characterizes cliques among
connected diamond-free sources and among connected complete multipartite
sources. A fixed five-spin matrix does the same for joins of a nonempty
clique with disjoint nonempty cliques. Positive rational perturbations
exist for each source, but may depend on it; for triangle-free sources
a single strictly positive three-spin matrix suffices.
Further obstructions use degree-class adjacency and triangle counts.
Exact rational certificates cover all $989$ connected noncliques of
orders three through seven. The maximum-degree theorem has $821$
exact checks, including $819$ direct partition-coefficient comparisons.
An explicit ten-vertex graph proves that the small-parameter criterion
alone is incomplete, while the full hierarchy detects it.
These results are partial progress toward the connected-source
conjecture of Lee--Oh--Seo; the conjecture and completeness of the
full hierarchy remain unproved.
\end{abstract}

\section{Problem and scope}
For a finite simple graph $G=(V,\mathcal{E})$ and an entrywise nonnegative
symmetric $q\times q$ matrix $H$, write
\begin{equation}
 Z_G(\mathbf x;H)
 =\sum_{\sigma:V\to\{0,\ldots,q-1\}}
       \prod_{uv\in\mathcal{E}}H_{\sigma(u),\sigma(v)}
       \prod_{v\in V}x_{\sigma(v)}.                 \label{eq:Z}
\end{equation}
It is homogeneous of degree $n=|V|$. We call $H$
\emph{antiferromagnetic} when it has at most one positive eigenvalue,
counting multiplicity. This is the spectral convention used in
\cite{LOS,LS}.

Lee--Oh--Seo prove that $Z_G$ is Lorentzian for every antiferromagnetic
$H$ when $G$ is a clique \cite[Theorem 3.1]{LOS}. They conjecture a
converse; the connected-source formulation is stated explicitly in
\cite[Section 6]{LS}. We use that formulation:
\begin{quote}
If $G$ is connected and $Z_G(\mathbf x;H)$ is Lorentzian for every
antiferromagnetic $H$, must $G$ be a clique?
\end{quote}
The restriction to connected graphs matters, since partition functions
factor over components and Lorentzian polynomials are closed under
products. The present paper establishes obstructions for several classes
of source graphs, using only matrices of rank two.

These questions concern all antiferromagnetic targets, rather than
only proper-colouring targets. Matherne--Morales--Selover establish
Lorentzianity of chromatic symmetric functions for abelian Dyck paths
\cite[Theorem 1.5]{MMS}. Lee--Oh--Seo already give a non-Lorentzian
three-vertex-path example with target
$\left(\begin{smallmatrix}0&0&1\\0&0&2\\1&2&0\end{smallmatrix}\right)$,
as well as a $K_{3,3}$ example \cite[Section 7]{LOS}. We do not claim
those small obstructions as new. A research note by Zeng
\cite{Zeng} gives a fixed rank-three, three-spin obstruction for
bistable bipartite graphs with a two-vertex cut. Our triangle-free
result uses a strictly positive rank-two target and requires neither
bistability nor a two-vertex cut; it does not claim Zeng's coefficient
formula relating the obstruction to vertex connectivity.
Our proposed contributions are a maximum-degree criterion excluding
true twins at that degree (Theorem~\ref{thm:maxtwins}), the resulting
obstructions for all connected noncomplete edge-regular graphs,
and fixed targets for triangle-free, connected diamond-free and complete
multipartite sources,
a fixed five-spin target for joins of a clique with disjoint cliques,
the neighbourhood-independence-gap criterion, the degree-count
identity and its compression certificate, and the parameterized
conditioned-set identities below. A diamond means an induced copy
of $K_4$ with one edge removed. The four-spin target for this class
has one zero entry; the positive perturbations are allowed to depend
on the source graph.

The arguments are elementary counting and real symmetric linear
algebra. Their computational corroboration uses exact arithmetic.
A targeted prior-art search does not certify worldwide priority; a
separate referee assessment of this draft remains pending.

\section{A rank-two family and the Hessian test}
We use a necessary consequence of the definition of Lorentzian
polynomials \cite{BH}. If a homogeneous polynomial $f$ of degree
$n\geq2$ with nonnegative coefficients is Lorentzian, then
\begin{equation}
 \Hess(\dd^\alpha f)\text{ has at most one positive eigenvalue}
 \quad\text{for every }\alpha\in\mathbb N_0^q,\quad|\alpha|=n-2.
                                                        \label{eq:test}
\end{equation}
The derivative is quadratic, so its Hessian is constant. We need
only this necessary condition for the general obstructions.
For the finite fixed-target limits reported below, we also use the
full characterization: nonnegative homogeneous polynomials with
M-convex support are Lorentzian precisely when all these quadratic
derivatives satisfy the Hessian condition \cite[Theorem 2.25]{BH}.

\begin{lemma}\label{lem:target}
For a nonconstant vector $\mathbf t\in(0,\infty)^q$, the matrix
\begin{equation}
 H(\mathbf t)_{ij}=t_i+t_j                              \label{eq:H}
\end{equation}
is positive entrywise and has inertia $(1,1,q-2)$, where inertia lists
positive, negative and zero eigenvalues.
\end{lemma}
\begin{proof}
Write $H=\one\mathbf t^\top+\mathbf t\one^\top$.
Its rank is at most two, and its quadratic form vanishes on
$\one^\perp$. A positive definite two-dimensional subspace would
intersect $\one^\perp$ nontrivially, which is impossible.
For $t_i\ne t_j$, the corresponding principal determinant is
$4t_it_j-(t_i+t_j)^2=-(t_i-t_j)^2<0$. Thus there are a positive and a
negative eigenvalue, and rank is exactly two.
\end{proof}

For integer $r\geq0$, introduce
\begin{equation}
 z_i=\frac{1+t_i}{2},\qquad
 k_i=\frac{t_i-1}{t_i+1},\qquad
 k_i z_i^r=z_i^r-z_i^{r-1}\quad(r\geq1).
                                                        \label{eq:zk}
\end{equation}
A baseline spin of value $t=1$ has self-interaction $2$.
If two free vertices are adjacent, the correction to their separate
baseline changes is
\begin{equation}
 \frac{2(t_i+t_j)}{(1+t_i)(1+t_j)}=1-k_i k_j.
                                                        \label{eq:correction}
\end{equation}
All denominators in these identities are positive.

\section{A fixed target for all triangle-free sources}
It is useful to normalize the interaction with the baseline spin.
For distinct $a_i\in(-1,1)$, set
\begin{equation}
 K(\mathbf a)=\one\one^\top-\mathbf a\mathbf a^\top,
 \qquad K_{ij}=1-a_i a_j.                              \label{eq:K}
\end{equation}
This matrix is positive entrywise and has rank two and exactly one
positive eigenvalue. Indeed, its quadratic form is nonpositive on
$\one^\perp$, and a principal $2\times2$ determinant for $a_i\ne a_j$
equals $-(a_i-a_j)^2<0$.

For a graph with $n\geq3$ vertices, $m$ edges, degrees $d_v$ and
$T$ triangles, abbreviate
\begin{equation}
 N=n(n-1)(n-2),\quad E=2m(n-2),\quad
 W=\sum_v d_v(d_v-1),\quad U=6T.                      \label{eq:counts}
\end{equation}
Here $W$ counts ordered wedges, including those inside triangles.
Use $\mathbf a=(\ell,0,1/2)$, with $-1<\ell<0$, and let
$V_{i\cdot}=(1,a_i,a_i^2)$.

\begin{theorem}[Edge--wedge--triangle identity]\label{thm:normalized}
For every finite simple graph $G$ on $n\geq3$ vertices,
\begin{align}
 \Hess\bigl(\dd_0\dd_1^{\,n-3}Z_G(\mathbf x;K(\mathbf a))\bigr)
    &=(n-3)!\,V\widehat{B}_\ell V^\top,                 \label{eq:normalized}\\
 \widehat{B}_\ell&=
 \begin{pmatrix}
 N&-\ell E&0\\-\ell E&\ell^2W-E&\ell W\\
 0&\ell W&-\ell^2U
 \end{pmatrix},                                       \label{eq:Bhat}\\
 \det\widehat{B}_\ell
    &=\ell^2\bigl[N(EU-W^2)+\ell^2U(E^2-NW)\bigr].    \label{eq:detnormalized}
\end{align}
If $W^2>EU$, a positive rank-two three-spin target makes $Z_G$
non-Lorentzian.
\end{theorem}
\begin{proof}
Place the conditioned spin $0$ at $v$ and baseline spin $1$, with
$a_1=0$, everywhere else. Every initial edge weight is one.
After assigning free spins $i,j$ to distinct vertices $u,w\ne v$,
the exact edge product is
\[
 (1-\ell a_i A_{uv})(1-\ell a_j A_{vw})
                         (1-a_i a_j A_{uw}).
\]
Summing over ordered distinct triples gives counts $N$ for the
constant term, $E$ for every single-edge indicator, $W$ for each
two-edge product, and $U$ for the three-edge product. This is exactly
$V\widehat{B}_\ell V^\top$. The derivative multiplicity is $(n-3)!$.
Expanding the $3\times3$ determinant proves
\eqref{eq:detnormalized}.

If $A_*=N(W^2-EU)>0$ and $R_*=U(E^2-NW)$, choose
$L=2+\lceil\max(R_*,0)/A_*\rceil$ and $\ell=-1/L$.
Then $\ell^2R_*<A_*$, so the determinant is negative.
The entry $(\widehat{B}_\ell)_{00}=N$ is positive, and hence its
inertia is $(2,1,0)$. The Vandermonde matrix $V$ is invertible,
so \eqref{eq:test} proves the conclusion.
\end{proof}

\begin{corollary}[One fixed positive target]\label{cor:trianglefree}
Let
\begin{equation}
 K_* =\frac14\begin{pmatrix}3&4&5\\4&4&4\\5&4&3\end{pmatrix}.
                                                               \label{eq:Kstar}
\end{equation}
For a triangle-free finite simple graph $G$, $Z_G(\mathbf x;K_*)$
is Lorentzian if and only if every component of $G$ is an isolated
vertex or a single edge. In particular, every connected triangle-free
graph on at least three vertices fails the universal-target property.
\end{corollary}
\begin{proof}
The target is $K(-1/2,0,1/2)$ and has eigenvalues $3,-1/2,0$.
If $G$ contains two adjacent edges, then $n\geq3$, $W>0$ and $U=0$.
Thus \eqref{eq:detnormalized} becomes $-NW^2/4<0$, proving the
obstruction for this fixed target.
Otherwise every component is a vertex or an edge. Their partition
functions are respectively a nonnegative linear form and a
quadratic form with Hessian $2K_*$. They are Lorentzian, as is
their product \cite[Corollary 2.32]{BH}. The empty graph, if allowed,
contributes the constant one.
\end{proof}

For a regular graph, $E=nd(n-2)$, $W=nd(d-1)$ and $U=n\tau$;
the inequality $W^2>EU$ is exactly \eqref{eq:threshold} below.
The general inequality also applies to irregular graphs with triangles.

\section{The degree-count kernel}
Let $m=|\mathcal{E}|$, let $d_v$ denote vertex degree, and let
$\Delta=\max_v d_v\geq2$. For $0\leq r\leq\Delta$, let
$\nu_r=|\{v:d_v=r\}|$, and put $\bnu=(\nu_0,\ldots,\nu_\Delta)^\top$.
For $1\leq r,s\leq\Delta$, let
\begin{equation}
 C_{rs}=|\{(u,v):uv\in\mathcal{E},\ d_u=r,\ d_v=s\}|.
                                                        \label{eq:C}
\end{equation}
These are \emph{ordered} edge counts; in particular an edge inside a
degree class contributes two to its diagonal entry.
Let $e_0,\ldots,e_\Delta$ be the standard basis, and let
$L$ be the $(\Delta+1)\times\Delta$ matrix with column
$L_r=e_r-e_{r-1}$ for $1\leq r\leq\Delta$.
Define the symmetric integer matrix
\begin{equation}
 B_G=\bnu\bnu^\top-\diag(\bnu)-LCL^\top.                \label{eq:B}
\end{equation}

\begin{theorem}[Degree-count identity]\label{thm:kernel}
Take $q=\Delta+1$, $t_i=i+1$ for $0\leq i\leq\Delta$, and
$H=H(\mathbf t)$. With $V_{ir}=z_i^r$, one has
\begin{equation}
 \Hess\bigl(\dd_0^{\,n-2}Z_G(\mathbf x;H)\bigr)
 =(n-2)!\,2^m\,V B_G V^\top.                          \label{eq:kernel}
\end{equation}
The matrix $V$ is invertible. Consequently, two positive eigenvalues
of $B_G$ certify that $Z_G(\mathbf x;H)$ is not Lorentzian.
\end{theorem}
\begin{proof}
The $(i,j)$ entry of the Hessian is $(n-2)!$ times a sum over ordered
distinct vertices $(u,v)$, giving them spins $i,j$ and every other
vertex the baseline spin $0$, whose value is $t_0=1$.
This interpretation is valid even if a free spin equals the baseline
or the two free spins agree: ordered placements give precisely the
falling-factorial multiplicities of differentiation.

The baseline edge product is $2^m$. Changing the two free vertices
separately gives $z_i^{d_u}z_j^{d_v}$; if $uv\in\mathcal{E}$,
\eqref{eq:correction} supplies the correction. After dividing by
$(n-2)!2^m$, the Hessian entry is therefore
\[
 \sum_{u\ne v}z_i^{d_u}z_j^{d_v}
 -\sum_{(u,v):uv\in\mathcal{E}} k_i z_i^{d_u} k_j z_j^{d_v}.
\]
The first sum is the $(i,j)$ entry of
$V(\bnu\bnu^\top-\diag(\bnu))V^\top$.
Using \eqref{eq:zk}, the second is the corresponding entry of
$VLCL^\top V^\top$. This proves \eqref{eq:kernel}.
The $z_i=(i+2)/2$ are distinct, so the square Vandermonde matrix $V$
is invertible. Sylvester's law of inertia and \eqref{eq:test}
complete the proof.
\end{proof}

\begin{theorem}[Independent minimum-degree class]\label{thm:minimum}
Suppose $G$ has positive minimum degree $\delta$ and its
minimum-degree vertices form an independent set. Then $B_G$ has at
least two positive eigenvalues. In particular,
$Z_G(\mathbf x;H(1,2,\ldots,\Delta+1))$ is not Lorentzian.
\end{theorem}
\begin{proof}
Put $p=\one\in\R^{\Delta+1}$ and $u=e_{\delta-1}$.
Since $L^\top p=0$, $\nu_{\delta-1}=0$ and $C_{\delta\delta}=0$,
\begin{equation}
 a:=p^\top B_Gp=n(n-1)>0,\qquad
 u^\top B_Gp=0,\qquad u^\top B_Gu=0.                   \label{eq:null}
\end{equation}
In computing $L^\top u$, a possible component indexed $\delta-1$
is harmless: that degree class is empty. Hence
\[
 B_Gu=LCe_\delta\ne0,
\]
where on the right $e_\delta$ denotes the column basis vector in
$\R^\Delta$ indexed by the degree $\delta$.
For the inequality, a minimum-degree vertex has a neighbour, and
every such neighbour has degree greater than $\delta$. Choose the
largest $s$ for which $C_{s\delta}>0$; the coefficient of $e_s$ in
$LCe_\delta$ is then positive.

Set $w=B_Gu$ and $c=u^\top B_Gw=\|B_Gu\|_2^2>0$.
The compression to the columns $(p,u,w)$ has the form
\begin{equation}
 Q=
 \begin{pmatrix}
   a&0&b\\0&0&c\\b&c&e
 \end{pmatrix},
 \qquad \det Q=-ac^2<0.                              \label{eq:compression}
\end{equation}
Its positive diagonal entry $a$ excludes three negative eigenvalues,
so its inertia is $(2,1,0)$. In particular its three columns are
independent, and its two-dimensional positive subspace gives one for
$B_G$. Apply Theorem~\ref{thm:kernel}.
\end{proof}

\begin{corollary}\label{cor:leaf}
Every connected graph on at least three vertices with a leaf has a
positive rank-two antiferromagnetic target for which its partition
function is not Lorentzian. This holds for every tree of order at
least three.
\end{corollary}
\begin{proof}
Two degree-one vertices cannot be adjacent in a connected graph of
order at least three. Thus its minimum-degree class is independent.
Every tree of that order has a leaf.
\end{proof}

The certificate \eqref{eq:compression} uses integer matrices of size
at most $\Delta+1$, followed by a $3\times3$ determinant. It does not
require expanding the $q^n$ terms in \eqref{eq:Z}.

\section{A parameterized obstruction for regular graphs}
The preceding criterion is not sufficient for regular graphs.
Indeed, for a $d$-regular graph its kernel is
\begin{equation}
 B_G=n(n-1)e_de_d^\top
       -nd(e_d-e_{d-1})(e_d-e_{d-1})^\top,             \label{eq:regularfirst}
\end{equation}
which has at most one positive eigenvalue.
We therefore use a different quadratic derivative.

Let $G$ be $d$-regular, with $n\geq3$ and $d\geq2$, and let $T$
be its number of unordered triangles. Set $\tau=6T/n$.
Choose $s>0$, $s\notin\{1,2\}$, and take three spin values
\[
 (t_0,t_1,t_2)=(s,1,2),\qquad
 \ell=\frac{s-1}{s+1}.
\]
Define $z_i,k_i$ by \eqref{eq:zk}, and set
\[
 V_{i\cdot}=(1,k_i,k_i^2),\qquad
 D_z=\diag(z_0^d,z_1^d,z_2^d).
\]
The three $k_i$ are distinct, and $D_z$ is positive and invertible.

\begin{theorem}[Regular-graph identity]\label{thm:regular}
For the target $H=H(s,1,2)$,
\begin{align}
 &\Hess\bigl(\dd_0\dd_1^{\,n-3}Z_G(\mathbf x;H)\bigr)
       =c_G(s)D_z V B_\ell V^\top D_z,               \label{eq:regular}\\
 &c_G(s)=(n-3)!\,n\,2^{m-d}(s+1)^d>0, \notag\\
 &B_\ell=
 \begin{pmatrix}
 (n-1)(n-2)&-\ell d(n-2)&0\\
 -\ell d(n-2)&\ell^2d(d-1)-d(n-2)&\ell d(d-1)\\
 0&\ell d(d-1)&-\ell^2\tau
 \end{pmatrix}.                                     \label{eq:Bell}
\end{align}
Its determinant is
\begin{equation}
 \det B_\ell=\ell^2d(n-2)
 \left[(n-1)\bigl((n-2)\tau-d(d-1)^2\bigr)
              +\ell^2\tau(n-1-d)\right].             \label{eq:det}
\end{equation}
\end{theorem}
\begin{proof}
The Hessian entry is $(n-3)!$ times a sum over ordered distinct
vertices $(v,u,w)$: assign spin $0$ to $v$, free spins $i,j$ to
$u,w$, and baseline spin $1$ to all remaining vertices.
The initial edge weight, before changing $u,w$, is
$2^{m-d}(s+1)^d$, independent of $v$.
For a free vertex adjacent to $v$, its change has an additional factor
\[
 \frac{2(s+t_i)}{(s+1)(1+t_i)}=1-\ell k_i.
\]
The correction for adjacency of the two free vertices is still
\eqref{eq:correction}. Divide out
$(n-3)!2^{m-d}(s+1)^d z_i^d z_j^d$.
The remaining entry is
\begin{equation}
 \sum_{\substack{v,u,w\\\text{all distinct}}}
 (1-\ell k_i A_{uv})(1-\ell k_j A_{vw})
                         (1-k_i k_j A_{uw}),         \label{eq:triples}
\end{equation}
where $A$ is the adjacency matrix of $G$.

In this sum there are $n(n-1)(n-2)$ ordered triples.
Each single-edge indicator has sum $nd(n-2)$.
Each product of two different edge indicators has sum $nd(d-1)$:
choose its common vertex, then two ordered distinct neighbours.
The product of all three indicators has sum $6T$.
Expanding \eqref{eq:triples} and dividing by $n$ gives
\[
 \begin{split}
 &(n-1)(n-2)-\ell d(n-2)(k_i+k_j)\\
 &\quad+\bigl(\ell^2d(d-1)-d(n-2)\bigr)k_i k_j\\
 &\quad+\ell d(d-1)(k_i^2k_j+k_ik_j^2)
                         -\ell^2\tau k_i^2k_j^2,
 \end{split}
\]
which is $V_{i\cdot}B_\ell V_{j\cdot}^{\top}$.
This proves \eqref{eq:regular}. Expanding the determinant of
\eqref{eq:Bell} gives \eqref{eq:det}.
\end{proof}

\begin{corollary}[Triangle threshold]\label{cor:regular}
If
\begin{equation}
 \frac{6T}{n}<\frac{d(d-1)^2}{n-2},                  \label{eq:threshold}
\end{equation}
then a positive three-spin rank-two antiferromagnetic target makes
$Z_G$ non-Lorentzian. Conversely, if \eqref{eq:threshold} fails,
none of the Hessians in Theorem~\ref{thm:regular} has two positive
eigenvalues. This converse concerns only that derivative and target
family.
\end{corollary}
\begin{proof}
Write
\[
 A_*=(n-1)\bigl(d(d-1)^2-(n-2)\tau\bigr),\qquad
 R_*=\tau(n-1-d)\geq0.
\]
Under \eqref{eq:threshold}, $A_*>0$. The entirely rational choice
\begin{equation}
 K=2+\left\lceil R_*/A_*\right\rceil,\qquad
 \ell=-1/K,\qquad s=(K-1)/(K+1)                     \label{eq:choice}
\end{equation}
has $0<s<1$ and $\ell^2R_*<A_*$. Indeed, for $r=R_*/A_*\geq0$,
$(2+\lceil r\rceil)^2>r$.
Equation~\eqref{eq:det} is therefore negative.
Since $(B_\ell)_{00}>0$, its inertia is $(2,1,0)$.
The invertible congruence \eqref{eq:regular}, followed by
\eqref{eq:test}, proves the obstruction.

If \eqref{eq:threshold} fails, then $A_*\leq0$, $\tau>0$,
and \eqref{eq:det} is nonnegative for every allowed $\ell$.
The diagonal entries $(B_\ell)_{00}>0$ and
$(B_\ell)_{22}=-\ell^2\tau<0$ force at least one positive and one
negative eigenvalue. A nonnegative determinant then forces the third
to be nonpositive. The same holds for the congruent Hessian.
\end{proof}

\begin{corollary}
Every regular triangle-free graph of degree at least two has a
non-Lorentzian partition function for the fixed positive target
$H(1/3,1,2)$.
\end{corollary}
\begin{proof}
For $T=0$, \eqref{eq:det} is negative for every $\ell\ne0$.
Take $\ell=-1/2$, equivalently $s=1/3$.
\end{proof}

\begin{remark}
The strict threshold is exact for this one-conditioned-spin family,
not an optimal condition among all targets or derivatives. For a
clique $d=n-1$ and $\tau=d(d-1)$, the bracket in
\eqref{eq:det} vanishes identically, consistent with the known
Lorentzian result.
\end{remark}

\section{Conditioning larger vertex sets}
The single conditioned vertex can be replaced by a set. This yields
an explicit hierarchy of tests; no assertion is made that the
hierarchy detects every connected nonclique.

Fix any finite simple graph $G$, an integer $1\leq h\leq n-2$, and
$-1<\ell<0$. Take $q=h+2$ distinct values $a_i\in(-1,1)$ with
$a_0=\ell$, $a_1=0$, and the positive target $K(\mathbf a)$.
For $S\subseteq V$, $|S|=h$, write
\[
 r_S(v)=|N(v)\cap S|\quad(v\notin S),\qquad
 \nu^S_r=|\{v\notin S:r_S(v)=r\}|\quad(0\leq r\leq h).
\]
Let $C^S_{ab}$ count ordered edges of $G[V\setminus S]$ between
vertices of profiles $a,b$, and let $e(S)$ be the number of edges
inside $S$. Define the $(h+2)\times(h+1)$ coefficient matrix
\[
 F_{jr}=\binom rj(-\ell)^j\quad
 (0\leq j\leq h+1,\ 0\leq r\leq h),
\]
where $\binom rj=0$ for $j>r$.
Let $J$ be the $(h+2)\times(h+2)$ matrix that shifts coordinates
down by one, with zeros in the first row. Set
\begin{equation}
 \begin{split}
 \mathcal B_h(\ell)=\sum_{|S|=h}(1-\ell^2)^{e(S)}
 \bigl[&F\bigl(\boldsymbol\nu^S(\boldsymbol\nu^S)^\top
                    -\diag(\boldsymbol\nu^S)\bigr)F^\top\\
       &-(JF)C^S(JF)^\top\bigr].
 \end{split}                                           \label{eq:hierarchyB}
\end{equation}
All matrices in the brackets have size $h+2$.

\begin{theorem}[Conditioned-set identity]\label{thm:hierarchy}
With $V_{ij}=a_i^j$ for $0\leq i,j\leq h+1$, one has
\begin{align}
 \Hess\bigl(\dd_0^h\dd_1^{\,n-h-2}Z_G(\mathbf x;K(\mathbf a))\bigr)
    &=c_h V\mathcal B_h(\ell)V^\top,
                                                        \label{eq:hierarchy}\\
 c_h&=h!(n-h-2)!>0. \notag
\end{align}
Thus $\mathcal B_h(\ell)$ and the quadratic-derivative Hessian have
the same inertia. For $h=1$, $\mathcal B_1(\ell)=\widehat{B}_\ell$;
for a $d$-regular graph this is $nB_\ell$.
\end{theorem}
\begin{proof}
Assign spin $0$ to $S$ and baseline spin $1$ to its complement.
All edges outside $G[S]$ have weight one, and its $e(S)$ internal
edges have weight $1-\ell^2$. The baseline product is therefore
$(1-\ell^2)^{e(S)}$.
Changing a free vertex $u\notin S$ to spin $i$ multiplies its
separate baseline contribution by
$(1-\ell a_i)^{r_S(u)}$.
For two adjacent free vertices the additional correcting factor
is $1-a_i a_j$, as before.

The coefficient vector of $(1-\ell k)^r$ in powers of $k$ is the
$r$th column of $F$; the vector for $k(1-\ell k)^r$ is the
corresponding column of $JF$.
Summing over distinct free vertices first gives the first matrix
in \eqref{eq:hierarchyB}, and subtracting the adjacency correction
gives the second. Summing the baseline weights over $S$ gives
$\mathcal B_h(\ell)$.
The derivative multiplicity is $h!(n-h-2)!$: it counts ordered
placements of the $h$ conditioned spins and the remaining baseline
spins. This also accounts for coincidences between free and
conditioned spin labels. These observations prove
\eqref{eq:hierarchy}.
The $a_i$ are distinct, so the congruence is invertible.
For $h=1$, expanding the profiles gives precisely the triple counts
in Theorem~\ref{thm:normalized}.
\end{proof}

For fixed $h$, the kernel can be assembled by examining the
$\binom nh$ conditioned sets and their degree profiles; no expansion
over all $(h+2)^n$ spin assignments is needed.
The formulas extend continuously to $\ell=-1$, where
$(1-\ell^2)^{e(S)}$ retains exactly the independent sets $S$
(with exponent zero interpreted as weight one).

For comparison with the additive targets, \eqref{eq:zk} gives
$H(\mathbf t)=2\diag(\mathbf z)K(\mathbf k)\diag(\mathbf z)$.
When $G$ is $d$-regular,
\begin{equation}
 Z_G(\mathbf x;H(\mathbf t))
    =2^m Z_G((z_i^d x_i)_i;K(\mathbf k)).              \label{eq:regularscaling}
\end{equation}
Thus the additive-target version of \eqref{eq:hierarchy} has
the additional positive factor $2^{m-hd}(s+1)^{hd}$ and the
diagonal congruence $D_z$, when $t_0=s$ and $t_1=1$.
Such a source-independent diagonal change is not asserted for
irregular graphs.


\section{Structural consequences of the profile identity}
A change of polynomial basis makes the boundary of the hierarchy
particularly useful. Retain the notation $r_S$, $\bnu^S$ and $C^S$
from Theorem~\ref{thm:hierarchy}. Let $\overline{\bnu}^S$ append a
zero to $\bnu^S$, and let $D_h$ be the $(h+2)\times(h+1)$ matrix
with column $r$ equal to $e_r-e_{r+1}$, for $0\leq r\leq h$.
For $0\leq\theta\leq1$ define
\begin{equation}
\begin{split}
 Q_h(\theta)=\sum_{|S|=h}(1-\theta)^{e(S)}
 \bigl[&\theta\bigl(\overline{\bnu}^S(\overline{\bnu}^S)^\top
                -\diag(\overline{\bnu}^S)\bigr)\\
       &-D_h C^S D_h^\top\bigr].
\end{split}                                               \label{eq:profilekernel}
\end{equation}
At $\theta=1$, the convention $0^0=1$ leaves precisely the
independent $h$-sets.

\begin{lemma}[Profile-basis identity]\label{lem:profile}
Let $\theta=\ell^2$, with $-1<\ell<0$. For any $q\geq2$,
take $q$ distinct parameters in $(-1,1)$ with $a_0=\ell$, $a_1=0$,
and the target $K(\mathbf a)$. Put $U_{ir}=(1-\ell a_i)^r$ for
$0\leq i<q$ and $0\leq r\leq h+1$. Then
\begin{equation}
 \Hess\bigl(\dd_0^h\dd_1^{n-h-2}Z_G(\mathbf x;K(\mathbf a))\bigr)
  =\frac{h!(n-h-2)!}{\theta}\,UQ_h(\theta)U^\top.       \label{eq:profileidentity}
\end{equation}
The matrix $U$ is invertible when $q=h+2$; for other $q$ the
identity is a rectangular matrix factorization. It extends to $\ell=-1$
when $a_0=-1$ and the remaining $a_i$ stay distinct in $(-1,1)$.
\end{lemma}
\begin{proof}
In the polynomial basis $1,y,\ldots,y^{h+1}$, where
$y=1-\ell k$, the profile factor is $y^r$, and
\[
 k(1-\ell k)^r=\frac{y^r-y^{r+1}}{\ell}.
\]
Thus the adjacency correction in the proof of
Theorem~\ref{thm:hierarchy} has coefficient matrix
$D_hC^SD_h^\top/\ell^2$. The other term is the embedded profile-count
matrix in \eqref{eq:profilekernel}. This proves
\eqref{eq:profileidentity}, with the same derivative multiplicity.
The counting argument does not require the evaluation matrix to
be square. The nodes $1-\ell a_i$ are distinct, so $U$ is an
invertible Vandermonde matrix when $q=h+2$. All coefficients are
continuous at $\ell=-1$, proving the boundary statement; in the
square case invertibility also persists there.
\end{proof}

\subsection{The small-parameter limit and true twins}
The other endpoint of the profile kernel gives a different structural
criterion. Define the unweighted edge-profile matrix
\begin{equation}
 T_h=\sum_{|S|=h}C^S,\qquad 1\leq h\leq n-2.
 \label{eq:unweightedprofile}
\end{equation}
The sum here includes all $h$-sets, whether or not they are independent.

\begin{lemma}[A small-parameter obstruction]\label{lem:smallparameter}
If $T_h$ has a negative eigenvalue, a strictly positive rational
antiferromagnetic target of rank two on $h+2$ spins makes $Z_G$
non-Lorentzian.
\end{lemma}
\begin{proof}
Choose $f$ with $f^\top T_hf<0$. The matrix $D_h$ has full column
rank, so there is $y\in\R^{h+2}$ with $D_h^\top y=f$. Let
$p=\one\in\R^{h+2}$ and $R=n-h\geq2$. Polynomial continuation
of \eqref{eq:profilekernel} to zero gives
\[
 Q_h(0)=-D_hT_hD_h^\top,\qquad Q_h(0)p=0,\qquad
 d_0:=y^\top Q_h(0)y=-f^\top T_hf>0.
\]
For $\theta\geq0$, direct evaluation gives
\begin{align}
 p^\top Q_h(\theta)p&=\theta A(\theta),&
 A(\theta)&=R(R-1)\sum_{|S|=h}(1-\theta)^{e(S)},\notag\\
 p^\top Q_h(\theta)y&=\theta B(\theta),&
 B(\theta)&=(R-1)\sum_{|S|=h}(1-\theta)^{e(S)}
                         (\overline{\bnu}^S)^\top y.
 \label{eq:smallcompression}
\end{align}
In particular $A(0)=\binom nhR(R-1)>0$. The determinant of the
compression of $Q_h(\theta)$ to columns $p,y$ is
\[
 \theta A(\theta)y^\top Q_h(\theta)y-\theta^2 B(\theta)^2
   =\theta A(0)d_0+O(\theta^2).
\]
For all sufficiently small $\theta>0$, this determinant and the first
diagonal entry are positive. Hence the compression is positive definite,
and $Q_h(\theta)$ has at least two positive eigenvalues.

Choose $\ell=-1/m$ for a sufficiently large integer $m$, so that
$\theta=1/m^2$, and put $q=h+2$. Take $a_0=\ell$, $a_1=0$ and,
for example, $a_i=(i-1)/(q+1)$ for $2\leq i<q$. These rational
parameters are distinct and lie in $(-1,1)$. For each such nonzero
$\ell$, the matrix $U$ in Lemma~\ref{lem:profile} is invertible.
Thus the Hessian in \eqref{eq:profileidentity} has two positive
eigenvalues, and \eqref{eq:test} proves the claim. The argument does
not assume that $U$ remains invertible at $\ell=0$.
\end{proof}

Two distinct vertices are \emph{true twins} if their closed
neighbourhoods are equal. Such vertices are adjacent and have equal degree.

\begin{theorem}[Maximum degree and true twins]\label{thm:maxtwins}
Let $G$ be a connected noncomplete graph with maximum degree $\Delta$.
If no two vertices of degree $\Delta$ are true twins, a strictly positive
rational antiferromagnetic target of rank two on $\Delta+1$ spins has
a non-Lorentzian partition polynomial on $G$.
\end{theorem}
\begin{proof}
Put $h=\Delta-1$. Connectedness and noncompleteness imply
$\Delta\geq2$, and hence $1\leq h\leq n-2$.
If an ordered edge $uv$ outside an $h$-set $S$ contributes to
$(T_h)_{hh}$, both endpoints are adjacent to every vertex of $S$
and to one another. The degree bound forces
\[
 d_u=d_v=\Delta,\qquad N(u)\setminus\{v\}=S=N(v)\setminus\{u\}.
\]
Then $N[u]=N[v]$, contrary to the hypothesis. Therefore $(T_h)_{hh}=0$.
Take a vertex $u$ of degree $\Delta$ and a neighbour $v$, and use
$S=N(u)\setminus\{v\}$. The edge $uv$ contributes to $(T_h)_{hj}$
with $j=|N(v)\cap S|<h$, by the same argument. Thus $(T_h)_{hj}>0$.
The principal matrix on indices $h,j$ has determinant
$-(T_h)_{hj}^2<0$, so $T_h$ has a negative eigenvalue.
Apply Lemma~\ref{lem:smallparameter}, with $h+2=\Delta+1$.
\end{proof}

This proof is constructive. For $c=(T_h)_{hj}>0$ and $a=(T_h)_{jj}$,
the rational vector $f=e_j-(a+1)e_h/(2c)$ satisfies $f^\top T_hf=-1$.
Set $y_{h+1}=0$ and $y_r=\sum_{k=r}^h f_k$, so that $D_h^\top y=f$.
Testing $\ell=-2^{-k}$ with exact arithmetic eventually gives a
positive definite compression in \eqref{eq:smallcompression}.

\begin{corollary}\label{cor:edgeregular}
The conclusion of Theorem~\ref{thm:maxtwins} holds for every connected
noncomplete graph without true twins, for every connected graph of
order at least three with a unique maximum-degree vertex, and for every
connected noncomplete edge-regular graph. In particular it holds for
all connected noncomplete strongly regular graphs.
\end{corollary}
\begin{proof}
The first two statements follow directly. For the third, edge-regular
means regular of degree $d$, with a constant number $\lambda$ of
common neighbours for each adjacent pair. A true-twin edge would give
$\lambda=d-1$, making every edge a true-twin edge. Equality of closed
neighbourhoods along paths would then make the connected graph complete.
Thus no such edge exists. Strong regularity implies edge regularity.
\end{proof}

For an infinite family outside the earlier triangle threshold, take the
line graph of $K_m$, $m\geq6$. It has $n=m(m-1)/2$, $d=2(m-2)$ and
$\lambda=m-2$; its ordered triangle count per vertex is $\tau=d\lambda$.
Since
\[
 \lambda(n-2)-(d-1)^2
     =\frac{(m-3)(m^2-8m+14)}2>0,
\]
the threshold \eqref{eq:threshold} fails, while
Corollary~\ref{cor:edgeregular} still gives an obstruction.
Any nonclique counterexample to the full connected-source conjecture
must therefore contain true twins of maximum degree. Neither the number
of spins nor the parameter choice in this theorem is asserted uniform
over all graphs.

\subsection{The independent-set endpoint}

For a graph with an edge, let $\alpha$ denote independence number,
with $\alpha(\varnothing)=0$, and set
\begin{align}
 b(G)&=\max_{(u,v):uv\in\mathcal E}
             \alpha\bigl(G[N(u)\setminus\{v\}]\bigr),\\
 c(G)&=\max_{uv\in\mathcal E}
             \alpha\bigl(G[N(u)\cap N(v)]\bigr).       \label{eq:gapparameters}
\end{align}
The first maximum uses both orientations of each edge.

\begin{theorem}[Neighbourhood independence gap]\label{thm:gap}
If $b(G)>c(G)$, a strictly positive rational antiferromagnetic
matrix of rank two on $c(G)+3$ spins has a non-Lorentzian partition
function on $G$.
\end{theorem}
\begin{proof}
Put $h=c(G)+1$. Since $b(G)\geq h$, some edge $uv$ and independent
$h$-set $S\subseteq N(u)\setminus\{v\}$ exist. In particular,
$h\leq n-2$. Let
\[
 C_* = \sum_{\substack{|S|=h\\S\text{ independent}}}C^S,
 \qquad Q=Q_h(1).
\]
One has $(C_*)_{hh}=0$. Otherwise the two endpoints of an edge
would be adjacent to every vertex in an independent $h$-set,
contradicting $c(G)<h$. The chosen $u,v,S$ gives
$(C_*)_{hj}>0$ for some $j<h$, so $C_*e_h\ne0$.

Take $p=\one\in\R^{h+2}$ and $u_0=e_{h+1}$. Since
$D_h^\top p=0$, $D_h^\top u_0=-e_h$, and the last entry of
$\overline{\bnu}^S$ is zero,
\begin{align}
 s_0:=p^\top Qp&=I_h(G)(n-h)(n-h-1)>0,\\
 p^\top Qu_0&=0,\qquad u_0^\top Qu_0=0,\qquad
 Qu_0=D_hC_*e_h\ne0,                                  \label{eq:gapnull}
\end{align}
where $I_h(G)$ counts independent $h$-sets. The last inequality
uses injectivity of $D_h$.
Put $w=Qu_0$ and $c_0=\|Qu_0\|_2^2>0$. The compression of $Q$
to the columns $(p,u_0,w)$ has the form
\[
 \begin{pmatrix}s_0&0&b_0\\0&0&c_0\\b_0&c_0&e_0\end{pmatrix},
 \qquad \det=-s_0c_0^2<0.
\]
Its positive diagonal entry excludes three negative eigenvalues,
so the compression has two positive eigenvalues. Hence so does $Q$.
Choose any distinct rational remaining $a_i\in(-1,1)$ with $a_1=0$.
Lemma~\ref{lem:profile} first gives an obstruction at $\ell=-1$.
The two positive eigenvalues persist when the rational parameter
$\ell>-1$ is sufficiently close to $-1$. All $a_i$ then lie in
$(-1,1)$, so $K(\mathbf a)$ is strictly positive and has rank two.
The Hessian test \eqref{eq:test} proves the result.
\end{proof}

The proof is constructive: with the remaining rational $a_i$ fixed,
exactly testing $\ell=-1+1/k$ for integers $k\geq2$ eventually
finds a positive target. It does not assert a uniform bound on $k$
over all source graphs.

\begin{corollary}[One target for connected diamond-free sources]
\label{cor:diamond}
Let
\begin{equation}
 K_{\diamond}=\frac19
 \begin{pmatrix}
 0&9&12&15\\9&9&9&9\\12&9&8&7\\15&9&7&5
 \end{pmatrix}.                                      \label{eq:diamondtarget}
\end{equation}
For every connected diamond-free finite simple graph $G$,
$Z_G(\mathbf x;K_{\diamond})$ is Lorentzian if and only if
$G$ is a clique. Thus the connected-source characterization
conjecture holds on diamond-free graphs. For each nonclique,
a strictly positive rational rank-two target on at most four
spins also gives an obstruction.
\end{corollary}
\begin{proof}
The displayed target is $K(-1,0,1/3,2/3)$; it is nonnegative,
has rank two and inertia $(1,1,2)$. For a clique the result is
\cite[Theorem 3.1]{LOS}.

Suppose $G$ is not complete. If $G$ is triangle-free, then
$c(G)=0$ and $b(G)\geq1$, because a connected nonclique contains
two adjacent edges. The proof of Theorem~\ref{thm:gap}, at
$h=1$ and $\ell=-1$, gives a two-positive-eigenvalue Hessian for
the principal target $K(-1,0,1/3)$. This Hessian is a principal
submatrix of the same quadratic derivative for $K_{\diamond}$,
so the latter also fails \eqref{eq:test}.

Otherwise $G$ contains a triangle. Diamond-freeness implies
$c(G)=1$: a common neighbourhood of an edge cannot contain an
independent pair, while a triangle gives a common neighbour.
There is a nonsimplicial vertex $u$ of degree at least three.
Indeed, if every vertex of degree at least three were simplicial
and such a vertex $v$ existed, its closed neighbourhood $N[v]$
would be a clique of size at least four. Every vertex of that clique
has degree at least three and is simplicial. An outside neighbour
of any of them would therefore also be adjacent to $v$, a
contradiction. Hence $N[v]$ would be a whole component, making
$G$ complete. If no vertex has degree at least three, a connected
graph containing a triangle is just $K_3$, again complete.

Choose nonadjacent $x,y\in N(u)$ and a third neighbour $v$.
Then $\{x,y\}\subseteq N(u)\setminus\{v\}$, so $b(G)\geq2>c(G)$.
The proof of Theorem~\ref{thm:gap}, now at $h=2$ and $\ell=-1$,
applies directly to $K_{\diamond}$. Finally perturb its first
parameter from $-1$ into $(-1,0)$; the two positive Hessian
eigenvalues persist and all target entries become positive.
\end{proof}

The fixed matrix \eqref{eq:diamondtarget} has one zero entry.
Its strictly positive perturbation may depend on $G$; no fixed
positive four-spin target for the entire class is asserted.
The gap condition in Theorem~\ref{thm:gap} is sufficient, not
necessary for an obstruction or for failure of the universal-target
property.

\section{Clique separators with disjoint complete branches}
The independent-set boundary also treats sources containing diamonds
that need not satisfy the neighbourhood-gap criterion. Write $\vee$
for graph join and $\sqcup$ for disjoint union.

\begin{proposition}[Two branches]\label{prop:twocliques}
For integers $a,b,s\geq1$, let $G=K_s\vee(K_a\sqcup K_b)$.
Then $Z_G(\mathbf x;K_{\diamond})$ is not Lorentzian.
More generally, the same conclusion holds for any four-spin target
$K(-1,0,c,d)$ with $-1<c,d<1$ and all four parameters distinct.
\end{proposition}
\begin{proof}
Put $A=a-1$, $B=b-1$, $L=A+B$ and $C=A(A-1)+B(B-1)$.
Suppose first that $L>0$ and use $h=2$, $\theta=1$ in
Lemma~\ref{lem:profile}. Every independent two-set contains one
vertex in each branch, so there are $ab$ such sets. Outside each
one the profile counts are $\nu_1=L$, $\nu_2=s$, and the ordered
edge counts are $C_{11}=C$, $C_{12}=C_{21}=sL$ and
$C_{22}=s(s-1)$. Consequently,
\begin{equation}
 Q_2(1)=ab(0\oplus M),\qquad
 M=\begin{pmatrix}
 2AB&C&sL\\ C&2sL-C&s(s-1-L)\\
 sL&s(s-1-L)&-s(s-1)
 \end{pmatrix}.                                      \label{eq:twoM}
\end{equation}
Direct expansion gives
\begin{align}
 \det M&=-s(L+s-1)
 \bigl[L^2(L+2s-1)+2AB(s-1)(L+s)\bigr]<0,             \label{eq:twoMdet}\\
 \one^\top M\one&=(L+s)(L+s-1)>0.\notag
\end{align}
Thus $M$ has inertia $(2,1,0)$. The square evaluation matrix $U$
is invertible for the four distinct parameters, so the Hessian in
\eqref{eq:profileidentity} has two positive eigenvalues.
Here $n-h=L+s\geq2$, as required.

If $L=0$, then $a=b=1$ and $G$ is $K_{s+2}$ with one edge removed.
The counts in \eqref{eq:counts} are
\[
 N=s(s+1)(s+2),\quad E=s^2(s+3),\quad
 W=s(s^2+3s-2),\quad U=s(s+4)(s-1).
\]
They satisfy $W^2-EU=E^2-NW=4s^2$. The polynomial counting
identity \eqref{eq:normalized} extends to $\ell=-1$, and its
coefficient matrix does not depend on the third spin parameter.
Therefore, at $h=1$,
\[
 \det\widehat B_{-1}
 =N(EU-W^2)+U(E^2-NW)=-24s^3<0.
\]
Its first diagonal entry $N$ is positive. The three distinct
parameters $(-1,0,c)$ give an invertible evaluation matrix, hence
a principal Hessian submatrix with two positive eigenvalues.
The four-spin Hessian has at least as many positive directions.
In both cases \eqref{eq:test} proves the conclusion.
\end{proof}

\begin{theorem}[One fixed five-spin target]\label{thm:clusterjoin}
Let
\begin{equation}
 K_5=\frac1{16}
 \begin{pmatrix}
 0&16&20&24&28\\
 16&16&16&16&16\\
 20&16&15&14&13\\
 24&16&14&12&10\\
 28&16&13&10&7
 \end{pmatrix}
 =K(-1,0,1/4,1/2,3/4).                              \label{eq:fixedfive}
\end{equation}
For every $s\geq1$, $r\geq2$ and $a_1,\ldots,a_r\geq1$, the
connected noncomplete graph
\[
 G=K_s\vee\bigl(K_{a_1}\sqcup\cdots\sqcup K_{a_r}\bigr)
\]
has non-Lorentzian partition polynomial $Z_G(\mathbf x;K_5)$.
The target is nonnegative of rank two, with exactly one zero entry.
For each such $G$, a strictly positive rational five-spin
perturbation also gives an obstruction.
\end{theorem}
\begin{proof}
The target has inertia $(1,1,3)$ by the argument for
\eqref{eq:K}. For $r=2$, Proposition~\ref{prop:twocliques}
applies to the first four spins; its quadratic-derivative Hessian
is a principal submatrix of the five-spin Hessian.
We therefore assume $r\geq3$.

Put $A_i=a_i-1$, $L=\sum_i A_i$, and
$C=\sum_i A_i(A_i-1)$. If $L>0$, use $h=r$ at the boundary
$\theta=1$ in Lemma~\ref{lem:profile}, now with $q=5$.
Every independent $h$-set chooses one vertex from each branch,
and there are $P=\prod_i a_i$ choices. The remaining branch
vertices have profile $1$, and separator vertices have profile $h$.
The only nonzero part of $Q_h(1)$ is on indices $(1,2,h,h+1)$,
where it is $P$ times
\begin{equation}
 M=\begin{pmatrix}
 L(L-1)-C&C&0&sL\\
 C&-C&sL&-sL\\
 0&sL&0&s(s-1)\\
 sL&-sL&s(s-1)&-s(s-1)
 \end{pmatrix}.                                     \label{eq:clusterM}
\end{equation}
Its determinant and the principal determinant on profile indices
$2,h$ are
\begin{align}
 \det M&=Ls^2(L+s-1)\bigl[sL^2-(s-1)C\bigr]>0,
                                                      \label{eq:clusterMdet}\\
 \det M[\{2,h\}]&=-s^2L^2<0.\notag
\end{align}
Indeed $0\leq C\leq L(L-1)$, so the bracket in
\eqref{eq:clusterMdet} is at least $L(L+s-1)>0$.
The negative principal determinant gives a positive and a negative
direction. The positive full determinant then forces inertia
$(2,2,0)$.

The four active columns of $U$ have degrees $1,2,h,h+1$.
On the four positive nodes $1,5/4,3/2,7/4$ they are linearly
independent. To see this, a nonzero polynomial with four distinct
powers has at most three positive roots by Descartes' sign rule;
dependence would give four distinct positive roots. Thus the
$5\times4$ active evaluation matrix has full column rank.
Multiplication on its two sides preserves the nonzero inertia of
$M$, and \eqref{eq:profileidentity} gives two positive eigenvalues.
The derivative is valid because $n-h=L+s\geq2$.

If $L=0$, all branches are singletons. Choose $h=r-1\geq2$.
Each independent $h$-set leaves one branch vertex of profile $0$
and $s$ vertices of profile $h$. There are $r$ such sets, and
$Q_h(1)$ is $r$ times the following matrix on indices $(0,1,h,h+1)$:
\begin{equation}
 M_0=\begin{pmatrix}
 0&0&0&s\\0&0&s&-s\\
 0&s&0&s(s-1)\\s&-s&s(s-1)&-s(s-1)
 \end{pmatrix}.                                      \label{eq:clusterMzero}
\end{equation}
Here $\det M_0=s^4>0$, whereas its principal determinant on
indices $1,h$ is $-s^2$. Hence its inertia is again $(2,2,0)$.
The active powers $0,1,h,h+1$ are distinct, and the same four-node
argument gives full column rank. The Hessian therefore has two
positive eigenvalues. Now $n-h=s+1\geq2$.

Finally, perturb the first parameter from $-1$ into $(-1,0)$.
The two strictly positive Hessian eigenvalues persist by continuity.
Choose a rational perturbation; all target entries become strictly
positive, with rank two and one positive eigenvalue.
\end{proof}

When $r=1$ the source is a clique and its partition polynomial
is Lorentzian by \cite[Theorem 3.1]{LOS}. Thus the fixed target
\eqref{eq:fixedfive} characterizes the complete members of the
whole family with $s\geq1$, $r\geq1$ and nonempty branches.
The number of target states is independent of the number and sizes
of the branches. We do not assert that five is minimal, that one
strictly positive perturbation works for every source, or that the
four-spin target of Proposition~\ref{prop:twocliques} suffices for
the entire family.



\begin{remark}[A limit of a four-spin restriction]\label{rem:fourrestriction}
Full column rank in the preceding proof is essential. Take
$G=K_s\vee(K_2\sqcup K_1\sqcup K_1)$ and $h=3$. For
$K_{\diamond}$ the first row and column of the boundary Hessian
vanish. After removing its positive scalar factor, the remaining
matrix is $VMV^\top$, where $M$ is \eqref{eq:clusterM} with
$L=1,C=0$ and $V$ evaluates powers $1,2,3,4$ at $1,4/3,5/3$.
Its determinant is
\[
 \det(VMV^\top)=\frac{12800}{4782969}s^3(185s-473)>0
 \qquad(s\geq3).
\]
The original $M$ has two positive and two negative eigenvalues,
so the restriction has at most two positive eigenvalues. Its
positive determinant and positive trace therefore give inertia
$(1,2,0)$. This derivative is inconclusive for $s\geq3$.
Other derivatives may still detect these graphs; this is not a
counterexample to universality of $K_{\diamond}$.
\end{remark}

\section{Complete multipartite sources}
The same four-spin target also detects a family not contained in the
diamond-free class. In this section a complete multipartite graph has
nonempty parts, and two vertices are adjacent exactly when they belong
to different parts.

\begin{theorem}[One fixed target for complete multipartite sources]
\label{thm:multipartite}
For every connected complete multipartite graph $G$, the partition
polynomial $Z_G(\mathbf x;K_{\diamond})$ is Lorentzian if and only
if $G$ is a clique. For every noncomplete such $G$, a strictly positive
rational four-spin rank-two perturbation of $K_{\diamond}$ also has
a non-Lorentzian partition polynomial. The perturbation may depend on $G$.
\end{theorem}
\begin{proof}
The clique direction follows from \cite[Theorem 3.1]{LOS}.
For the converse, write the part sizes as $a_1,\ldots,a_r$.
Connectedness and noncompleteness give $r\geq2$ and
$\alpha=\max_i a_i\geq2$. Put $h=\alpha-1$; then
$1\leq h\leq n-2$. Apply Lemma~\ref{lem:profile} at
$\ell=-1$, $\theta=1$, and $q=4$. Only independent $h$-sets contribute.

Such a set $S$ lies in one part $i$ with $a_i\geq h$.
There are $w_i=\binom{a_i}{h}$ choices. Outside $S$, the
$k_i=a_i-h\in\{0,1\}$ vertices in that part have profile $0$,
and the other $y_i=n-a_i$ vertices have profile $h$.
Every pair between these two classes is an edge, there are no edges
within profile $0$, and the ordered edge count within profile $h$ is
\[
 t_i=y_i^2-\sum_{j\ne i}a_j^2.
\]
Summing over all $i$ with $a_i\geq h$, define
\begin{equation}
 A=\sum_i w_i\sum_{j\ne i}a_j(a_j-1),\qquad
 B=\sum_i w_i k_i y_i,\qquad C=\sum_i w_i t_i.
 \label{eq:multipartiteABC}
\end{equation}
All three numbers are nonnegative. Moreover $B>0$: a largest
part has $k_i=1$ and $y_i>0$.

First let $h\geq2$. Substitution into \eqref{eq:profilekernel}
shows that $Q_h(1)$ is supported on indices $(0,1,h,h+1)$,
where it is the matrix
\begin{equation}
 M=\begin{pmatrix}
 0&0&0&B\\0&0&B&-B\\0&B&A&C\\B&-B&C&-C
 \end{pmatrix}.                                    \label{eq:multipartiteM}
\end{equation}
Indeed, the profile-$0$ diagonal is zero because $k_i(k_i-1)=0$,
and its pair term with profile $h$ cancels the corresponding edge
term. The profile-$h$ diagonal is $y_i(y_i-1)-t_i$, whose sum is $A$.
All other entries follow from the two differences in the columns of $D_h$.
Now $\det M=B^4>0$, while its principal determinant on indices
$1,h$ is $-B^2<0$. Thus $M$ has inertia $(2,2,0)$.

The active evaluation matrix in \eqref{eq:profileidentity} evaluates
the powers $0,1,h,h+1$ at $0,1,4/3,5/3$. It is invertible.
In fact, a vanishing linear combination has zero constant coefficient
by evaluation at zero. Dividing by the variable leaves at most three
nonzero monomials, which by Descartes' sign rule have at most two
positive roots unless the polynomial is zero. The three distinct
positive evaluation nodes therefore force all coefficients to vanish.
The quadratic-derivative Hessian has two positive eigenvalues.

If $h=1$, the middle two profile indices coincide. Direct substitution
instead gives the $3\times3$ kernel
\[
 M_1=\begin{pmatrix}
 0&0&B\\0&A+2B&C-B\\B&C-B&-C
 \end{pmatrix}.
\]
Its determinant is $-B^2(A+2B)<0$, and its middle diagonal entry
is positive, so its inertia is $(2,1,0)$. The evaluation matrix on
powers $0,1,2$ at the same four nodes has full column rank.
Again the Hessian has two positive eigenvalues.

Both cases violate \eqref{eq:test}. Finally, replace the first target
parameter $-1$ by a rational $\ell\in(-1,0)$ sufficiently close to
$-1$. The two positive eigenvalues persist by continuity. All target
entries become strictly positive, and its rank and inertia are those
of the rank-two family already proved above.
\end{proof}

The theorem includes arbitrarily many parts and arbitrarily large part
sizes. Its proof uses only four active profile columns even when
$h>2$; a square evaluation matrix of size $h+2$ is unnecessary.
The Lorentzian characterization of complete multipartite graphs by
Choudhury--Khare \cite[Theorem 2.3]{CK} concerns instead the homogenized
determinant of a distance-matrix blowup. It is not a statement about
the spin partition polynomial \eqref{eq:Z} or the fixed target here.


\section{Examples and limits}
\subsection{A second-level obstruction beyond the triangle threshold}
Start with two disjoint copies of $K_4$ on vertices $\{0,1,2,3\}$
and $\{4,5,6,7\}$. Delete edges $(0,1)$ and $(4,5)$, and add
$(0,4)$ and $(1,5)$. The resulting graph is connected, cubic,
and has four triangles. Hence $\tau=3$, while the right side of
\eqref{eq:threshold} is $2$. The one-conditioned-spin test cannot
detect it for any of its allowed parameter values.

For $h=2$ and the boundary value $\ell=-1$, direct evaluation gives
\[
 \mathcal B_2(-1)=
 \begin{pmatrix}
 480&480&60&0\\480&264&-120&-12\\
 60&-120&-148&-16\\0&-12&-16&-4
 \end{pmatrix}.
\]
Symmetric elimination without permutations has consecutive pivots
\[
 480,\quad -216,\quad -\frac{11}{2},\quad \frac{106}{33}.
\]
Its inertia is therefore $(2,2,0)$. Theorem~\ref{thm:hierarchy}
at the continuous boundary supplies a non-Lorentzian Hessian.
Its two positive eigenvalues persist for $\ell>-1$ sufficiently
close to $-1$, giving positive rank-two targets as well.
The supplement checks explicitly that $\ell=-99/101$ works, using
rational elimination of $\mathcal B_2(-99/101)$ and direct
comparison with the additively scaled target $H(1/100,1,2,3)$
via \eqref{eq:regularscaling}.

\subsection{Parameter variation enlarges the test}
Allowing the boundary value $s=0$ in the counting identity gives
$\ell=-1$ and the more restrictive sufficient condition
\begin{equation}
 \tau<\frac{(n-1)d(d-1)^2}{(n-1)^2-d}.               \label{eq:old}
\end{equation}
For a concrete strict improvement, take the graph on
$\{0,\ldots,11\}$ with edges
\[
\begin{split}
 &(0,1),(0,2),(0,8),(0,9),(1,4),(1,7),(1,9),(2,5),\\
 &(2,8),(2,11),(3,6),(3,7),(3,8),(3,11),(4,5),(4,7),\\
 &(4,10),(5,6),(5,9),(6,10),(6,11),(7,10),(8,9),(10,11).
\end{split}
\]
It is connected and $4$-regular, and has $7$ triangles.
Thus $n=12$, $d=4$, $\tau=7/2$.
At $s=0$, $\det B_{-1}=540>0$, so this Hessian has at most one
positive eigenvalue. At $s=2/3$, $\ell=-1/5$ and
\[
 \det B_{-1/5}=-\frac{2004}{125}<0.
\]
The positive target is
\[
 H(2/3,1,2)=
 \begin{pmatrix}4/3&5/3&8/3\\5/3&2&3\\8/3&3&4\end{pmatrix}.
\]
The supplement verifies the full Hessian identity for this graph
with rational arithmetic.

\subsection{The degree-count test does not settle all irregular graphs}
\label{ex:degreefailure}
Consider the connected six-vertex graph with edge set
\[
 \{(0,1),(0,4),(0,5),(1,2),(1,5),(2,3),(2,5),(3,4)\}.
\]
Its degree counts are $\nu_2=2$, $\nu_3=4$, and its kernel is
\[
 B_G=
 \begin{pmatrix}
 0&0&0&0\\0&-2&0&2\\0&0&-6&16\\0&2&16&2
 \end{pmatrix}.
\]
The characteristic polynomial is
$\lambda(\lambda^3+6\lambda^2-264\lambda-560)$.
The cubic has exactly one positive root by Descartes' rule of signs
and a sign change between zero and sufficiently large positive
$\lambda$. Thus this kernel has only one positive eigenvalue.
It does not establish Lorentzianity of the partition function.
A different derivative supplies a countermodel in the finite
certificates described below.

Graphs with many triangles can also fail
\eqref{eq:threshold}. The condition $b(G)>c(G)$
also leaves cases untreated. These sufficient criteria do not prove
the full connected-source conjecture, and their failure is not evidence
against that conjecture.


\subsection{Limits of three fixed subtargets}
The same fixed four-spin target $K_{\diamond}$ has an exact
quadratic-derivative obstruction for every connected nonclique of
orders three through seven in the graph atlas: $989$ graphs.
This is a stronger finite assertion than allowing a different
target for each graph, but it is not an arbitrary-order theorem.
An additional deterministic sample of $64$ graphs, with orders up
to ten, also yields exact obstructions. This sample is not an
exhaustive census beyond order seven.

Deleting one of the last three states can destroy this property.
For the principal target on states $(0,1,2)$, the seven-vertex graph
with edge set
\[
 \{(0,1),(0,4),(1,2),(1,4),(2,3),(4,5),(5,6)\}
\]
has a Lorentzian partition polynomial. For the principal targets
on states $(0,1,3)$ and $(0,2,3)$, the six-vertex graph in
Section~\ref{ex:degreefailure} has that property. The supplement
records every quadratic-derivative Hessian: $21$ in the first
case and $15$ in each of the other two. Exact rational congruences
give at most one positive eigenvalue in each, and the full M-convex
exchange condition is verified on each support. The full
characterization cited after \eqref{eq:test} proves these finite
claims. They concern three specific fixed targets; they do not
show that every three-spin target fails or that four spins are
necessary in general.

\subsection{The small-parameter criterion is not complete}
\label{ex:zerolimit}
Let $G=C_5[K_2]$: replace each vertex of a five-cycle by a two-vertex
clique, and join all vertices in two replacement cliques whenever the
original cycle vertices are adjacent. This is a connected noncomplete
$5$-regular graph on ten vertices, and each replacement pair is a pair
of true twins. Every matrix $T_h$, $1\leq h\leq8$, is positive
semidefinite, so Lemma~\ref{lem:smallparameter} does not detect $G$.
Here this is an exact finite claim with a short coefficient certificate.

For an ordered edge $uv$, put $c=|N(u)\cap N(v)|$,
$a=d_u-1-c$, $b=d_v-1-c$, and $d=n-d_u-d_v+c$.
Partitioning the other vertices by their adjacency to $u,v$ gives
the contribution
\begin{equation}
 \sum_k\binom ck\binom a{r-k}\binom b{t-k}
               \binom d{h-r-t+k}
 \label{eq:localprofile}
\end{equation}
to $(T_h)_{rt}$, with out-of-range binomial coefficients zero.
For this graph, ten ordered edges have $c=4$ and forty have $c=2$.
Thus $(T_h)_{rt}$ is the coefficient of $x^r y^t z^h$ in
\[
 10(1+xyz)^4(1+z)^4
 +40(1+xyz)^2(1+xz)^2(1+yz)^2(1+z)^2.
\]
For $h=1,2,3,4$, ordinary symmetric elimination without permutations
has the following positive diagonal pivots:
\begin{center}
\begin{tabular}{@{}cl@{}}
\toprule
$h$ & Pivots of $T_h$\\
\midrule
1 & $120,\ 200/3$\\
2 & $100,\ 224,\ 300/7$\\
3 & $40,\ 480,\ 360,\ 200/9$\\
4 & $10,\ 320,\ 600,\ 592/3,\ 250/37$\\
\bottomrule
\end{tabular}
\end{center}
For each fixed outside edge, complementing $S$ in its other eight
vertices replaces $h$ by $8-h$ and each profile $r$ by $4-r$.
It therefore relates $T_h$ and $T_{8-h}$ by reversal of the active
profile indices and zero padding. This proves positive semidefiniteness
for $h=5,6,7$;
for $h=8$ the sole nonzero entry is $(T_8)_{44}=50$.

This failure does not extend to the full profile hierarchy.
Use $q=h+2$, $a_0=\ell$, $a_1=0$ and $a_i=(i-1)/(h+2)$ for $i\geq2$.
At $\ell=-99/100$, its kernels for $h=3,4,5$ have respective inertias
$(2,3,0)$, $(2,4,0)$ and $(2,3,2)$, verified by exact congruences
and direct partition-coefficient identities. All three targets are
strictly positive. Since $\alpha(G)=2$, the unscaled endpoint $Q_h(1)$
vanishes at these levels. The first nonzero term after dividing by
the smallest power of $1-\theta$ has two positive eigenvalues for
$h=3,4$, but only one for $h=5$. The nearby kernel at $h=5$ nevertheless
has two. Thus neither the zero-parameter test nor the leading endpoint
term alone supplies a general completeness argument.


\section{Exact checks and reproducibility}
The accompanying supplement contains the source scripts, their data,
and a manifest of file hashes. All checks were run in the authoring
context; none is a formal proof-assistant verification or an
independent model assessment. The general proofs are the counting
identities and compression arguments above.

\begin{center}
\begin{tabular}{@{}p{9.0cm}r@{}}
\toprule
Check & Cases\\
\midrule
Normalized three-spin identity, including irregular sources &288\\
Degree identity, all connected graphs of orders $3$--$5$,
original nonnegative target & 29\\
Degree identity, strictly positive target of Theorem~\ref{thm:kernel}
&29\\
Independent minimum-degree compression,
connected graph atlas of orders $3$--$7$ &806\\
Original regular identity, exact enumeration &13\\
Parameterized regular identity: $18$ graphs, $5$ rational parameters
&90\\
Additive-target conditioned-set identity for $h=1,2,3$
&26\\
Normalized conditioned-set identity, including irregular sources &24\\
\bottomrule
\end{tabular}
\end{center}
\begin{center}
\begin{tabular}{@{}p{9.0cm}r@{}}
\toprule
Certificate census & Cases\\
\midrule
Exact three-spin countermodels for all connected noncliques
of orders $3$--$6$ in the atlas &137\\
Positive rank-two hierarchy certificates for all connected
noncliques of orders $3$--$7$ &989\\
Characteristic-polynomial audit of those kernels &989\\
Neighbourhood-gap compression and positive continuation &634\\
Fixed diamond-free target, direct identities &255\\
Positive perturbations of the diamond-free target, direct identities &255\\
Fixed four-spin certificates on all connected noncliques of orders $3$--$7$ &989\\
\bottomrule
\end{tabular}
\end{center}


\begin{center}
\begin{tabular}{@{}p{9.0cm}r@{}}
\toprule
Clique-separator and fixed-target checks & Cases\\
\midrule
Two-branch sources: direct fixed four-spin boundary identities &64\\
Two-branch sources: direct positive continuations &46\\
Clique--cluster joins: direct fixed five-spin boundary identities &64\\
Clique--cluster joins: direct positive continuations &27\\
Additional seeded fixed four-spin obstructions, orders up to ten &64\\
Exact Lorentzian counterexamples to three fixed subtargets &3\\
Four-spin restriction: direct checks of Remark~\ref{rem:fourrestriction} &6\\
Complete multipartite sources: direct fixed four-spin identities &36\\
Complete multipartite sources: direct positive continuations &17\\
\bottomrule
\end{tabular}
\end{center}

\begin{center}
\begin{tabular}{@{}p{9.0cm}r@{}}
\toprule
Small-parameter checks & Cases\\
\midrule
Maximum-degree theorem: exact compression certificates &821\\
Maximum-degree theorem: direct positive Hessian identities &819\\
Unweighted profile-kernel probe, atlas and selected blowups &1001\\
$C_5[K_2]$: exact semidefinite kernels, all valid $h$ &8\\
$C_5[K_2]$: direct positive full-hierarchy identities &3\\
\bottomrule
\end{tabular}
\end{center}

The maximum-degree checks cover all $817$ eligible connected noncliques
of orders three through seven, two larger wheels, the line graph of
$K_6$, and the $4$-by-$4$ rook graph. Each records a rational negative
direction for $T_h$ and a positive definite two-dimensional compression
of a nearby $Q_h$. The $15$- and $16$-vertex examples use exact kernels
and the proved identity; the other $819$ also use direct partition
enumeration. Among these $821$ graphs, $244$ fail the earlier
edge--wedge--triangle sufficient criterion. The separate unweighted
probe finds an indefinite $T_h$ for every atlas nonclique, but fails
on nine selected clique blowups beyond the atlas. Section~\ref{ex:zerolimit}
proves the first such limitation explicitly. Neither this probe nor
the theorem asserts completeness of the unweighted criterion.

The three-spin census contains $141$ connected graphs, including four
complete-graph controls. Each of the $137$ certificates records the
edge set, integer target $H_{ij}=t_i+t_j$, derivative multi-index,
integer Hessian, its negative determinant and positive trace.
These signs prove inertia $(2,1,0)$. The finite census does not
establish the conjecture at arbitrary order; absence of a detected
obstruction in the clique controls is not its proof.

The larger census contains all $994$ connected graphs of orders
three through seven, of which five are cliques. For each of the
$989$ noncliques it gives a strictly positive rational rank-two
target, an integer-scaled profile kernel, a rational congruence
certificate, and direct agreement with a Hessian obtained by
partition-coefficient enumeration. A separate audit checks the
characteristic polynomial of every stored kernel. Its real roots
are counted exactly using Descartes bounds for positive and negative
roots together with the exact zero multiplicity; the bounds sum to
the degree and agree with the congruence inertia. These are separate
computations within the same authoring context, not a separate referee.

The gap check covers $634$ graphs in the same atlas. All $255$
connected diamond-free noncliques there satisfy the criterion;
$88$ are triangle-free and $167$ contain triangles. For these
$255$ sources, both the fixed boundary target and strictly positive
rational perturbations were checked by direct partition-coefficient
enumeration. The finite checks corroborate the general proof of
Corollary~\ref{cor:diamond}; they do not replace it.

The clique-separator checks expand the determinant identities
symbolically and compare the predicted Hessians with direct
partition-coefficient enumeration. The fixed-five sample includes
$r=2,3,4,5,6$, so the number of conditioned vertices can exceed
the number of nonbaseline target states. The full proof uses the
rectangular identity and the four-column rank argument, rather
than extrapolating from these checks.
The complete multipartite sample has largest parts from two through six.
Its $36$ boundary identities and $17$ positive identities were checked
by direct partition-coefficient enumeration; the symbolic determinants
in Theorem~\ref{thm:multipartite} were expanded separately.

\noindent\begin{minipage}{\textwidth}
The main scripts are:
\begin{flushleft}\ttfamily
\begin{tabular}{@{}ll@{}}
verify\_degree\_kernel.py & verify\_neighbourhood\_gap.py\\
verify\_normalized\_target.py & verify\_fixed\_diamond.py\\
verify\_regular\_obstruction.py & verify\_two\_clique\_extension.py\\
verify\_parametric\_obstruction.py & verify\_fixed\_five\_cluster.py\\
verify\_conditioning\_hierarchy.py & probe\_fixed\_diamond\_scope.py\\
certify\_source\_witnesses.py & probe\_larger\_fixed\_target.py\\
certify\_profile\_hierarchy.py & probe\_three\_spin\_targets.py\\
audit\_profile\_certificates.py & verify\_four\_spin\_restriction.py\\
verify\_multipartite\_extension.py & verify\_maximum\_degree\_twins.py\\
verify\_blowup\_layer.py & probe\_zero\_parameter\_kernels.py\\
\end{tabular}
\end{flushleft}
\end{minipage}\par\medskip
The source-witness script reads the preserved finite-search candidates and
recomputes their Hessians by exact enumeration.
The parameterized script separately checks the symbolic determinant,
the rational parameter selection, and the positive-target identity.
Python, SymPy and NetworkX suffice for these checks; NumPy is used
only by the exploratory witness search.
The supplementary README gives invocation and provenance details.

\section{Remaining question}
The obstructions extract degree-to-degree adjacency and
edge, wedge and triangle counts. The conditioned-set hierarchy
uses finer adjacency profiles. Its small-parameter limit gives the
maximum-degree criterion: any nonclique counterexample to the source
conjecture must have true twins of maximum degree. The explicit
$C_5[K_2]$ limitation shows why that limit alone is insufficient. Its independent-set boundary gives
the neighbourhood-gap theorem and the diamond-free characterization.
Its rectangular evaluation also gives a fixed five-spin obstruction
for joins of a clique with any number of complete branches, and a
fixed four-spin obstruction for all noncomplete complete multipartite
sources.
A general argument would have to prove that a suitable test detects
every connected nonclique; no such sufficiency is established here.
The distinction between bounded
countermodel search and a proof for graphs of arbitrary order remains
essential.

\paragraph{Authorship and AI assistance.}
Lluis Eriksson is the responsible human author. Codex assisted with
literature searches, derivations, manuscript preparation and
computational checks. Any model assessments are maintained separately,
bound to the exact manuscript hash; they do not constitute human peer
review or formal verification. No numerical rating forms part of this
manuscript.
The author operates AIRR, the intended repository; this relationship
does not constitute external peer review.

\clearpage
\begin{thebibliography}{9}
\bibitem{BH}
P. Br\"and\'en and J. Huh,
\emph{Lorentzian polynomials},
Annals of Mathematics \textbf{192} (2020), 821--891.
\href{https://doi.org/10.4007/annals.2020.192.3.4}
{doi:10.4007/annals.2020.192.3.4}.

\bibitem{LOS}
J. Lee, J. Oh and J. Seo,
\emph{Counting homomorphisms in antiferromagnetic graphs via
Lorentzian polynomials}, arXiv:2506.13659v2 (2025).
\url{https://arxiv.org/html/2506.13659v2}.

\bibitem{LS}
J. Lee and J. Seo,
\emph{Antiferromagnetic models are clique-minimizing},
arXiv:2608.17920v1 (2026).
\url{https://arxiv.org/html/2608.17920v1}.

\bibitem{MMS}
J. P. Matherne, A. H. Morales and J. Selover,
\emph{The Newton polytope and Lorentzian property of chromatic
symmetric functions},
Selecta Mathematica \textbf{30} (2024), article 42;
arXiv:2201.07333v4.
\url{https://arxiv.org/html/2201.07333v4}.

\bibitem{Zeng}
X. Zeng,
\emph{Two-vertex cuts and a Lorentzian obstruction for bistable
bipartite graphs}, research note, AI-Has-Taste,
BW2-PR20260914-lorentzian\_cliques (2026).
\href{https://github.com/ArtificialZeng/AI-Has-Taste/blob/f3ae144b74f238f22c2b84b5283cb14ce912fd33/reproducible_source_code/BigWIN2/BW2-PR20260914-lorentzian_cliques/source/manuscript/main.tex}{Archived source, accessed 17 October 2026}.

\bibitem{CK}
P. N. Choudhury and A. Khare,
\emph{Blowup polynomials and delta-matroids of graphs},
S\'eminaire Lotharingien de Combinatoire \textbf{86B} (2022),
article 42, 12 pp.; arXiv:2203.04105.
\url{https://www.mat.univie.ac.at/~slc/wpapers/FPSAC2022/42.pdf}.

\end{thebibliography}
\end{document}
